\clearpage
\begin{summarybox}
    \textbf{\textcolor{CSummary}{一句话记住}}
    \textbf{闭区间上一直升或一直降，最大值和最小值看两端；
    图像还必须连续，才能把两端之间所有数写成完整值域。}

    \medskip
    \textbf{\textcolor{CSummary}{两个公式（闭区间、单调且连续）}}
    \[
        \begin{aligned}
            \text{递增：}\quad &\boxed{\text{值域}=[f(a),f(b)]},\\[3pt]
            \text{递减：}\quad &\boxed{\text{值域}=[f(b),f(a)]}.
        \end{aligned}
    \]
    \textbf{不想记顺序：}先算 $f(a),f(b)$，小的写左边，大的写右边。

    \medskip
    \textbf{\textcolor{CSummary}{三步快速检查}}
    \begin{itemize}[leftmargin=*,itemsep=0.28em,topsep=0.3em]
        \item \textbf{第一看范围：}定义域包含两个端点吗？区间内每一点都能代入吗？
        \item \textbf{第二看方向：}整段一直上升或下降吗？有没有对称轴、绝对值转折点？
        \item \textbf{第三看缺口：}图像中间连续吗？不能把跳过的数算进值域。
    \end{itemize}

    \medskip
    \textbf{\textcolor{CSummary}{证明方法速记}}
    \textbf{看大小关系}得到端点范围；\textbf{看连续图像或寻找零点}保证中间不缺值；
    对能解出 $x$ 的具体函数，也可\textbf{从任意 $y$ 倒着求出 $x$} 来核对。

    \medskip
    \textbf{\textcolor{CSummary}{一道题最容易漏的地方}}
    \textbf{“端点是最值”不等于“中间每个数都取到”！}
\end{summarybox}
